
% Body fragment for repository integration.
% Assumes standard AMS theorem environments and basic macros.
% A standalone wrapper is provided in casimir_divisor_core_transport_standalone.tex.

\chapter{Casimir Recurrence Modules, Divisor Cores, and Intrinsic Spectral Transport}
\label{chap:casimir-divisor}

\section{The first-principles problem}

The spectral discriminant
\[
\Delta_A(x)
\]
is the first non-scalar layer of the modular characteristic hierarchy. As long as it is treated
merely as a polynomial, one sees only its decategorified shadow: branch points, reciprocal
growth rates, and characteristic polynomials of recurrence matrices.
What is missing is an \emph{object} whose internal subquotients remember the divisor structure of
\(\Delta_A\), and whose functorial maps explain what Drinfeld--Sokolov reduction, embeddings,
and exact factorization functors can preserve.

\begin{quote}
The correct algebraic transport object attached to the \(L_2\)-layer is not the polynomial
\(\Delta_A\) but the cyclic \(k[t]\)-module cut out by its reversed polynomial.
Divisors of \(\Delta_A\) then become actual subquotients, not merely factors.
\end{quote}

This shift immediately forces a correction of language.
One does \emph{not} in general have a canonical projector onto a chosen spectral factor.
Projectors exist precisely on the coprime locus.
On the repeated-root locus, the correct invariant is a canonical
subquotient together with a canonical extension class.
That is where the Jordan geometry lives.

Throughout this chapter \(k\) denotes a field of characteristic \(0\).

\section{Pure algebra: cyclic transport objects and divisor cores}

\subsection{The universal cyclic transport object}

\begin{definition}[Reversal of a polynomial]
For a nonzero polynomial \(f(x)\in k[x]\) of degree \(d\), define its \emph{reversal}
by
\[
f^{\rev}(t):=t^d f(t^{-1}).
\]
Thus if \(f(x)=\prod_{i=1}^d (1-\lambda_i x)\) over a splitting field,
then
\[
f^{\rev}(t)=\prod_{i=1}^d (t-\lambda_i).
\]
\end{definition}

\begin{definition}[Universal cyclic transport object]
Let \(p(t)\in k[t]\) be monic of degree \(d\).
The \emph{cyclic transport object} attached to \(p\) is the \(k[t]\)-module
\[
M(p):=k[t]/(p).
\]
Multiplication by \(t\) defines an endomorphism
\[
T_p \colon M(p)\to M(p),\qquad T_p([f])=[tf].
\]
\end{definition}

\begin{proposition}[Universal property]
Let \(V\) be a finite-dimensional \(k\)-vector space, let \(T\in \End_k(V)\), and let \(v\in V\).
Assume that \(V\) is generated by \(v\) under \(T\), equivalently
\[
V = k[T]\cdot v,
\]
and that \(p(T)v=0\).
Then there is a unique surjective \(k[t]\)-module map
\[
\phi_v \colon M(p)\twoheadrightarrow V,\qquad [f(t)]\mapsto f(T)v.
\]
If \(p\) is the minimal polynomial of \(T\) on the cyclic vector \(v\), then \(\phi_v\) is an isomorphism.
\end{proposition}

\begin{proof}
The assignment \(f\mapsto f(T)v\) is \(k[t]\)-linear from \(k[t]\) to \(V\).
The relation \(p(T)v=0\) implies that \((p)\) lies in its kernel, so the map descends to
\[
M(p)=k[t]/(p)\to V.
\]
Surjectivity is exactly the cyclicity hypothesis \(V=k[T]\cdot v\).
If \(p\) is minimal on \(v\), then the kernel is precisely \((p)\), hence the induced map is injective.
\end{proof}

\begin{remark}
This is the minimal algebraic transport object.
Any cyclic finite-dimensional endomorphism system is obtained from it by quotient;
if the characteristic and minimal polynomials coincide, then the system is exactly \(M(p)\).
\end{remark}

\subsection{Divisor cores}

\begin{definition}[Divisor core]
Let \(p(t)\in k[t]\) be monic and let \(g(t)\mid p(t)\) be a monic divisor.
The \emph{divisor core} associated to \(g\) is the submodule
\[
\mathcal C_g(p):=\bigl(p/g\bigr)M(p)\subset M(p).
\]
\end{definition}

\begin{theorem}[Divisor-core calculus]
\label{thm:divisor-core-calculus}
Let \(p\in k[t]\) be monic, and let \(g,g_1,g_2\mid p\) be monic divisors.

\begin{enumerate}[label=\textup{(\roman*)}]
\item There is a canonical isomorphism
\[
\mathcal C_g(p)\cong k[t]/(g).
\]

\item There is a canonical short exact sequence
\[
0\longrightarrow \mathcal C_g(p)\longrightarrow M(p)\longrightarrow k[t]/(p/g)\longrightarrow 0.
\]

\item The divisor cores satisfy the lattice identities
\[
\mathcal C_{g_1}(p)\cap \mathcal C_{g_2}(p)=\mathcal C_{\gcd(g_1,g_2)}(p),
\]
\[
\mathcal C_{g_1}(p)+\mathcal C_{g_2}(p)=\mathcal C_{\operatorname{lcm}(g_1,g_2)}(p).
\]

\item The following are equivalent:
\begin{enumerate}[label=\textup{(\alph*)}]
\item \(\mathcal C_g(p)\) is the image of an idempotent polynomial in \(T_p\);
\item \(M(p)\cong \mathcal C_g(p)\oplus k[t]/(p/g)\);
\item \(\gcd(g,p/g)=1\).
\end{enumerate}
\end{enumerate}
\end{theorem}

\begin{proof}
For \textup{(i)}, define
\[
\Phi_g\colon k[t]\to \mathcal C_g(p),\qquad f\mapsto f\cdot (p/g)\bmod p.
\]
If \(f-f'=gh\), then
\[
(f-f')(p/g)=hp\equiv 0\pmod p,
\]
so \(\Phi_g\) descends to a surjective map
\[
k[t]/(g)\twoheadrightarrow \mathcal C_g(p).
\]
If \(\Phi_g([f])=0\), then \(p\mid f(p/g)\). Since \(p=(p/g)g\) and \(k[t]\) is a domain, this forces \(g\mid f\).
Hence the kernel is \((g)\), proving \(\mathcal C_g(p)\cong k[t]/(g)\).

For \textup{(ii)}, the quotient map
\[
M(p)=k[t]/(p)\longrightarrow k[t]/(p/g)
\]
has kernel consisting of classes divisible by \(p/g\), namely \(\mathcal C_g(p)\).

For \textup{(iii)}, write
\[
\mathcal C_{g_i}(p)=\frac{(p/g_i)}{(p)}.
\]
Intersection of submodules corresponds to intersection of the corresponding ideals modulo \((p)\):
\[
\frac{(p/g_1)}{(p)}\cap \frac{(p/g_2)}{(p)}
=\frac{(p/g_1)\cap (p/g_2)}{(p)}
=\frac{\bigl(\operatorname{lcm}(p/g_1,p/g_2)\bigr)}{(p)}
=\frac{(p/\gcd(g_1,g_2))}{(p)}.
\]
This is \(\mathcal C_{\gcd(g_1,g_2)}(p)\).
Similarly,
\[
\frac{(p/g_1)}{(p)}+\frac{(p/g_2)}{(p)}
=\frac{(p/g_1)+(p/g_2)}{(p)}
=\frac{\bigl(\gcd(p/g_1,p/g_2)\bigr)}{(p)}
=\frac{(p/\operatorname{lcm}(g_1,g_2))}{(p)}.
\]

For \textup{(iv)}, image of an idempotent is equivalent to direct summand decomposition.
By \textup{(i)} and \textup{(ii)}, this is equivalent to the Chinese remainder decomposition
\[
k[t]/(p)\cong k[t]/(g)\oplus k[t]/(p/g),
\]
which holds if and only if \((g)\) and \((p/g)\) are comaximal, i.e.\ if and only if \(\gcd(g,p/g)=1\).
\end{proof}

\begin{corollary}[Divisors classify submodules]
Every submodule of \(M(p)\) is of the form \(\mathcal C_g(p)\) for a unique monic divisor \(g\mid p\).
Thus the lattice of submodules of \(M(p)\) is isomorphic to the divisor lattice of \(p\).
\end{corollary}

\begin{proof}
Submodules of \(k[t]/(p)\) correspond to ideals \((a)\) with \((p)\subseteq (a)\subseteq k[t]\).
Because \(k[t]\) is a PID, such an ideal is generated by a unique monic divisor \(a\mid p\).
Writing \(a=p/g\), the corresponding submodule is \(\mathcal C_g(p)\).
\end{proof}

\begin{remark}
This is the first structural correction to the spectral story.
A factor \(g\mid p\) always gives a canonical submodule.
It gives a canonical projector \emph{only} when \(\gcd(g,p/g)=1\).
\end{remark}

\subsection{Maps are controlled by gcd}

\begin{theorem}[Hom \(=\) gcd]
\label{thm:hom-equals-gcd}
Let \(p,q\in k[t]\) be monic, and let \(g=\gcd(p,q)\).
Then there is a canonical isomorphism of \(k\)-vector spaces
\[
\Hom_{k[t]}(M(p),M(q))\cong k[t]/(g).
\]
In particular,
\[
\dim_k\Hom_{k[t]}(M(p),M(q))=\deg g.
\]
\end{theorem}

\begin{proof}
A \(k[t]\)-linear map \(f\colon M(p)\to M(q)\) is determined by the image of \(1\bmod p\).
Write
\[
x:=f(1)\in M(q)=k[t]/(q).
\]
The condition that \(f\) be \(k[t]\)-linear is exactly
\[
p\cdot x = f(p\cdot 1)=0 \quad \text{in } M(q),
\]
which means \(px\in (q)\).
Therefore
\[
\Hom_{k[t]}(M(p),M(q))
\cong \{x \in k[t]/(q): px\in (q)\}
\cong ((q):(p))/(q),
\]
where \((q):(p)\) is the colon ideal.

Write \(p=gp'\), \(q=gq'\) with \(\gcd(p',q')=1\).
Then
\[
px\in (q)
\iff gp'x\in (gq')
\iff q' \mid p'x.
\]
Because \(p'\) and \(q'\) are coprime, this is equivalent to \(q'\mid x\).
Hence
\[
((q):(p)) = (q').
\]
Therefore
\[
\Hom_{k[t]}(M(p),M(q))
\cong (q')/(gq')
\cong k[t]/(g).
\]
\end{proof}

\begin{definition}[Maximal common core]
Given monic \(p,q\in k[t]\), define
\[
g_{p,q}:=\gcd(p,q),
\qquad
\mathcal C^{\max}(p,q):=M(g_{p,q})=k[t]/(g_{p,q}).
\]
\end{definition}

\begin{theorem}[Universal factorization through the maximal common core]
\label{thm:factorization-through-common-core}
Let \(p,q\in k[t]\) be monic, let \(g=\gcd(p,q)\), and write
\[
\pi_p\colon M(p)\twoheadrightarrow M(g)
\]
for the canonical quotient, and
\[
\iota_q\colon M(g)\hookrightarrow M(q)
\]
for the canonical embedding onto the divisor core \(\mathcal C_g(q)\).
Then every \(k[t]\)-linear map
\[
f\colon M(p)\to M(q)
\]
factors uniquely as
\[
f=\iota_q\circ m_u\circ \pi_p
\]
for a unique class \(u\in M(g)\), where
\[
m_u\colon M(g)\to M(g)
\]
denotes multiplication by \(u\).
\end{theorem}

\begin{proof}
By Theorem~\ref{thm:hom-equals-gcd},
\[
\Hom_{k[t]}(M(p),M(q))\cong M(g).
\]
The map \(u\mapsto \iota_q\circ m_u\circ \pi_p\) is \(k\)-linear from \(M(g)\) into \(\Hom_{k[t]}(M(p),M(q))\).
It is injective because \(\pi_p\) is surjective and \(\iota_q\) is injective.
Since both source and target have dimension \(\deg g\), it is an isomorphism.
\end{proof}

\begin{corollary}[Spectral defects]
\label{cor:spectral-defects}
Let \(p,q\) be monic and \(g=\gcd(p,q)\).
Define the \emph{defect modules}
\[
\mathsf D_{p/q}:=k[t]/(p/g),
\qquad
\mathsf D_{q/p}:=k[t]/(q/g).
\]
Then there are canonical short exact sequences
\[
0\longrightarrow \mathcal C^{\max}(p,q)\longrightarrow M(p)\longrightarrow \mathsf D_{p/q}\longrightarrow 0,
\]
\[
0\longrightarrow \mathcal C^{\max}(p,q)\longrightarrow M(q)\longrightarrow \mathsf D_{q/p}\longrightarrow 0.
\]
\end{corollary}

\begin{proof}
Apply Theorem~\ref{thm:divisor-core-calculus}\textup{(ii)} with \(g=\gcd(p,q)\).
\end{proof}

\begin{remark}
Transport maps between cyclic spectral objects are never mysterious.
They are exactly multiplication operators on the maximal common core.
The common divisor controls everything; the two defects measure what each side carries that the other cannot see.
\end{remark}

\subsection{Primary and Jordan sectors}

\begin{definition}[Primary divisor cores]
Let \(p\in k[t]\) be monic, and let \(q\) be an irreducible factor of \(p\).
For \(1\le m\le \nu_q(p)\), define
\[
\mathcal J_q^m(p):=\mathcal C_{q^m}(p)\subset M(p).
\]
\end{definition}

\begin{theorem}[Primary Jordan filtration]
\label{thm:primary-jordan-filtration}
Let \(q\) be irreducible and \(n=\nu_q(p)\).
Then there is a canonical chain
\[
0\subset \mathcal J_q^1(p)\subset \mathcal J_q^2(p)\subset \cdots \subset \mathcal J_q^n(p)\subset M(p),
\]
and each successive quotient satisfies
\[
\mathcal J_q^m(p)/\mathcal J_q^{m-1}(p)\cong k[t]/(q).
\]
If \(q=t-\lambda\) splits linearly over \(k\), then
\[
\mathcal J_{t-\lambda}^m(p)\cong k[t]/((t-\lambda)^m),
\]
and \(T_p-\lambda\) acts on this sector by a single Jordan block of size \(m\).
\end{theorem}

\begin{proof}
By Theorem~\ref{thm:divisor-core-calculus}\textup{(i)},
\[
\mathcal J_q^m(p)\cong k[t]/(q^m).
\]
The inclusions follow from
\[
\frac{p}{q^m}=q\frac{p}{q^{m+1}}.
\]
The inclusion \(\mathcal J_q^{m-1}(p)\subset \mathcal J_q^m(p)\) corresponds,
under \(\mathcal J_q^m(p)\cong k[t]/(q^m)\), to multiplication by \(q\).
The quotient
\[
k[t]/(q^m)\Big/\; q\,k[t]/(q^m)\cong k[t]/(q)
\]
gives the successive quotients.

If \(q=t-\lambda\), write \(\tau=t-\lambda\).
Then
\[
k[t]/((t-\lambda)^m)\cong k[\tau]/(\tau^m),
\]
and multiplication by \(t\) is multiplication by \(\lambda+\tau\), so \(t-\lambda\) acts by the standard nilpotent Jordan shift.
\end{proof}

\begin{corollary}[Repeated roots are extension classes]
\label{cor:repeated-roots-extension}
Let \(g\mid p\), set \(h=p/g\), and let \(d_g=\gcd(g,h)\).
The canonical sequence
\[
0\longrightarrow \mathcal C_g(p)\longrightarrow M(p)\longrightarrow k[t]/(h)\longrightarrow 0
\]
defines an obstruction class
\[
\varepsilon_g(p)\in
\Ext^1_{k[t]}\!\bigl(k[t]/(h),k[t]/(g)\bigr)
\cong k[t]/(d_g).
\]
Under the standard resolution
\[
0\longrightarrow k[t]\xrightarrow{\;h\;}k[t]\longrightarrow k[t]/(h)\longrightarrow 0,
\]
the class \(\varepsilon_g(p)\) is represented by \(1\bmod d_g\).
Thus \(\varepsilon_g(p)=0\) if and only if \(\gcd(g,h)=1\).
If \(d_g\) is nonconstant, the divisor core is canonical but the extension is not split.
\end{corollary}

\begin{proof}
Identify \(\mathcal C_g(p)\) with \(k[t]/(g)\) by
Theorem~\ref{thm:divisor-core-calculus}\textup{(i)}.
The middle term \(M(p)=k[t]/(gh)\) is generated by the class \(e=1\bmod gh\);
the submodule is generated by \(n=h\bmod gh\).
The relations are
\[
g n=0,\qquad h e=n.
\]
The displayed resolution computes
\[
\Ext^1_{k[t]}(k[t]/(h),k[t]/(g))
\cong \operatorname{coker}\!\bigl(h:k[t]/(g)\to k[t]/(g)\bigr)
\cong k[t]/(g,h).
\]
In the presentation above the relation \(h e=n\) is exactly the class
\(1\bmod (g,h)\).
This class vanishes precisely when \(d_g=1\), equivalently when the
Chinese-remainder splitting of
Theorem~\ref{thm:divisor-core-calculus}\textup{(iv)} exists.
\end{proof}

\subsection{Squarefree spectral sheets}

\begin{definition}[Squarefree spectral sheet]
Let \(S(x)\in k[x]\) satisfy \(S(0)=1\) and suppose that over a splitting field \(K/k\),
\[
S(x)=\prod_{i=1}^m (1-\lambda_i x)
\]
with distinct \(\lambda_i\).
Define
\[
\operatorname{Sh}(S):=\bigoplus_{i=1}^m K e_i,
\qquad
T_S(e_i)=\lambda_i e_i.
\]
\end{definition}

\begin{theorem}[Squarefree spectral sheet theorem]
\label{thm:squarefree-sheet}
Let \(S(x)\in k[x]\) be squarefree with \(S(0)=1\), and let
\[
s(t):=S^{\rev}(t).
\]
Then over the splitting field \(K\),
\[
K\otimes_k k[t]/(s)\cong \operatorname{Sh}(S)
\]
as \(K[t]\)-modules.
Equivalently, multiplication by \(t\) on \(k[t]/(s)\) becomes semisimple with eigenvalues the reciprocals of the roots of \(S\).
\end{theorem}

\begin{proof}
Since
\[
s(t)=\prod_{i=1}^m (t-\lambda_i)
\]
with distinct roots, the Chinese remainder theorem gives
\[
K[t]/(s)\cong \bigoplus_{i=1}^m K[t]/(t-\lambda_i)\cong \bigoplus_{i=1}^m K e_i.
\]
On the \(i\)-th summand, multiplication by \(t\) is multiplication by \(\lambda_i\).
\end{proof}

\begin{proposition}[Coprime-locus functional calculus]
\label{prop:coprime-functional-calculus}
Let \(V\) be finite-dimensional over \(k\), let \(T\in \End_k(V)\) have characteristic polynomial \(p\), and let \(g\mid p\).
Assume \(\gcd(g,p/g)=1\).
Choose polynomials \(a,b\in k[t]\) such that
\[
a(t)g(t)+b(t)\frac{p(t)}{g(t)}=1.
\]
Then
\[
e_g(T):=b(T)\Bigl(\frac{p}{g}\Bigr)(T)
\]
is an idempotent, independent of the chosen Bezout pair modulo the relation above, and its image is a \(T\)-stable direct summand whose characteristic polynomial is \(g\).
\end{proposition}

\begin{proof}
The Bezout identity implies
\[
a(T)g(T)+b(T)\Bigl(\frac{p}{g}\Bigr)(T)=I_V.
\]
Since \(p(T)=0\), the action of \(k[t]\) on \(V\) factors through \(k[t]/(p)\).
The Chinese-remainder idempotent
\[
e_g=b(t)(p/g)\bmod p
\]
has residue \(1\) in \(k[t]/(g)\) and residue \(0\) in \(k[t]/(p/g)\).
Thus \(e_g(T)^2=e_g(T)\), and it is independent of the Bezout pair modulo \((p)\).
For the primary decomposition of \(V\) with respect to the coprime factors \(g\) and \(p/g\),
the image of \(e_g(T)\) is the summand whose characteristic polynomial is the
\(g\)-factor of the characteristic polynomial \(p\).
This is the same summand obtained from
\[
k[t]/(p)\cong k[t]/(g)\oplus k[t]/(p/g),
\]
functorially under the \(k[t]/(p)\)-module structure on \(V\).
\end{proof}

\begin{proposition}[No universal polynomial projector off the coprime locus]
\label{prop:no-projector}
Let \(p\in k[t]\) be monic and \(g\mid p\).
If \(\gcd(g,p/g)\neq 1\), then there is no polynomial \(e(t)\in k[t]\) whose action on \(M(p)\) is an idempotent with image \(\mathcal C_g(p)\).
\end{proposition}

\begin{proof}
If such an idempotent existed, then \(\mathcal C_g(p)\) would be a direct summand of \(M(p)\), contradicting Theorem~\ref{thm:divisor-core-calculus}\textup{(iv)}.
\end{proof}

\section{Casimir recurrence modules from the bar tower}

\subsection{The recurrence hypothesis}

Let \(A\) be a modular Koszul chiral algebra in the Lie-theoretic world.
Write
\[
Z_A:=k[C_1,\dots,C_r,K]
\]
for the commutative algebra generated by the rank-\(r\) Casimirs together with the level operator.

\begin{hypothesis}[Finite-character recurrence hypothesis]
\label{hyp:finite-character-recurrence}
Assume the following.

\begin{enumerate}[label=\textup{(\roman*)}]
\item Each \(H_n(\Bbar(A))\) has finite generalized \(Z_A\)-support.

\item The class sequence
\[
u_A(n):=[H_n(\Bbar(A))]\in K_0(Z_A\text{-mod}_{\mathrm{fl}})\otimes_{\mathbb Z} k
\]
satisfies a minimal monic shift recurrence
\[
p_A(\sigma)u_A=0
\]
for some monic polynomial \(p_A(t)\in k[t]\).

\item Writing
\[
\Delta_A(x)\in k[x],\qquad \Delta_A(0)=1,\qquad d_A:=\deg \Delta_A,
\]
the recurrence polynomial is the reverse polynomial
\[
p_A(t)=t^{d_A}\Delta_A(t^{-1}).
\]
\end{enumerate}
\end{hypothesis}

\begin{remark}
The preceding algebra does not depend on how \(p_A\) arises.
The hypothesis enters only to identify \(p_A\) with the discriminant of the bar generating function and to interpret the resulting module as the algebraic shadow of \(H^1_{\red}(A)\).
\end{remark}

\subsection{Intrinsic algebraic reduced cohomology}

\begin{definition}[Intrinsic Casimir recurrence module]
Under Hypothesis~\ref{hyp:finite-character-recurrence}, define
\[
\mathcal R_A:=k[\sigma]\cdot u_A
\subset \Bigl(K_0(Z_A\text{-mod}_{\mathrm{fl}})\otimes k\Bigr)^{\mathbb N}.
\]
Define the \emph{algebraic reduced first cohomology} by
\[
H^{1,\alg}_{\red}(A):=\mathcal R_A^\vee,
\]
and let
\[
KS_A:=\sigma^\vee\in \End\bigl(H^{1,\alg}_{\red}(A)\bigr).
\]
\end{definition}

\begin{theorem}[Minimal intrinsic realization]
\label{thm:minimal-intrinsic-realization}
Under Hypothesis~\ref{hyp:finite-character-recurrence},
\[
\mathcal R_A\cong k[t]/(p_A(t))
\]
as \(k[t]\)-modules, where \(t\) acts by \(\sigma\).
Consequently
\[
\dim_k H^{1,\alg}_{\red}(A)=\deg p_A=d_A,
\]
and
\[
\det(1-x\,KS_A)=\Delta_A(x).
\]
\end{theorem}

\begin{proof}
Consider the \(k[t]\)-linear map
\[
\Phi_A\colon k[t]\to \mathcal R_A,\qquad f(t)\mapsto f(\sigma)u_A.
\]
It is surjective by definition of \(\mathcal R_A\).
Its kernel is the annihilator ideal of \(u_A\), generated by the minimal monic annihilating polynomial \(p_A\).
Hence
\[
\mathcal R_A\cong k[t]/(p_A).
\]
Dualizing gives
\[
H^{1,\alg}_{\red}(A)\cong (k[t]/(p_A))^\vee.
\]

The characteristic polynomial of multiplication by \(t\) on \(k[t]/(p_A)\) is \(p_A\), so
\[
\det(1-x\,KS_A)=x^{d_A}p_A(x^{-1})=\Delta_A(x).
\]
\end{proof}

\begin{remark}
This construction is already intrinsic to the bar tower and the Casimir action.
It does \emph{not} require a prior closed formula for the generating function.
The only external input is the existence of the minimal recurrence polynomial.
\end{remark}

\subsection{A single Casimir can separate the spectrum}

\begin{lemma}[Generic Casimir separator]
\label{lem:generic-separator}
Let \(Z\) be a finitely generated commutative \(k\)-algebra over an infinite field \(k\), and let
\[
\chi_1,\dots,\chi_s\colon Z\to \bar{k}
\]
be pairwise distinct algebra homomorphisms.
Then there exists \(Q\in Z\) such that
\[
\chi_i(Q)\neq \chi_j(Q)\qquad\text{for all }i\neq j.
\]
In particular, for a finite set of joint generalized \(Z_A\)-characters, a generic linear combination of the Casimirs and level separates them.
\end{lemma}

\begin{proof}
For each pair \(i\neq j\), choose \(z_{ij}\in Z\) with \(\chi_i(z_{ij})\neq \chi_j(z_{ij})\).
Choose generators \(z_1,\dots,z_N\) of \(Z\) and write
\[
Q=\sum_{\nu=1}^N c_\nu z_\nu.
\]
For each pair \(i\neq j\), the condition \(\chi_i(Q)=\chi_j(Q)\) is a nontrivial linear equation in the coefficients \(c_\nu\).
Thus the bad choices of \((c_1,\dots,c_N)\in k^N\) lie in a finite union of proper hyperplanes.
Since \(k\) is infinite, their complement is nonempty.
\end{proof}

\begin{theorem}[Casimir moment reconstruction]
\label{thm:casimir-moment-reconstruction}
Assume there is an integer \(n_0\) such that, for every \(n\ge n_0\), the generalized
\(Z_A\)-spectrum occurring in \(H_n(\Bbar(A))\) is contained in a finite set
\[
\Xi_A=\{\chi_1,\dots,\chi_s\}.
\]
Choose \(Q\in Z_A\) separating \(\Xi_A\), and write
\[
q_i:=\chi_i(Q).
\]
For \(0\le m\le s-1\), define the scalar moment sequences
\[
\tau_m(n):=\operatorname{Tr}\!\bigl(Q^m\mid H_n(\Bbar(A))\bigr).
\]
Then for every \(n\ge n_0\), the multiplicity vector of generalized characters in \(H_n(\Bbar(A))\) is uniquely determined by
\[
\bigl(\tau_0(n),\dots,\tau_{s-1}(n)\bigr).
\]
Equivalently, the tail cyclic module generated by \((u_A(n))_{n\ge n_0}\) is recoverable from finitely many scalar Casimir traces.
\end{theorem}

\begin{proof}
Write
\[
H_n(\Bbar(A))\cong \bigoplus_{i=1}^s m_i(n)\,V_{\chi_i}
\]
in the Grothendieck group, where \(m_i(n)\in k\) denotes the multiplicity of the generalized \(\chi_i\)-block.

On the generalized \(\chi_i\)-block, the operator \(Q\) acts as \(q_i I + N_i\) with \(N_i\) nilpotent.
Hence
\[
\operatorname{Tr}(Q^m\mid V_{\chi_i})
=\operatorname{Tr}\bigl((q_iI+N_i)^m\bigr)
= q_i^m \dim V_{\chi_i},
\]
because every positive power of \(N_i\) has trace \(0\).
Therefore
\[
\tau_m(n)=\sum_{i=1}^s q_i^m\,m_i(n).
\]
Let \(m(n)=(m_i(n))_{i=1}^s\) and \(\tau(n)=(\tau_m(n))_{m=0}^{s-1}\).
Then
\[
\tau(n)=V\,m(n),
\]
where \(V=(q_i^m)\) is the Vandermonde matrix.
Since the \(q_i\) are pairwise distinct, \(V\) is invertible.
Thus \(m(n)=V^{-1}\tau(n)\).
\end{proof}

\begin{corollary}[Finite scalar probes suffice]
\label{cor:finite-scalar-probes}
Under the hypotheses of Theorem~\ref{thm:casimir-moment-reconstruction}, the tail cyclic \(k[t]\)-module generated by \((u_A(n))_{n\ge n_0}\) may be reconstructed from the finite vector-valued scalar sequence
\[
\tau_A(n):=\bigl(\tau_0(n),\dots,\tau_{s-1}(n)\bigr)\in k^s.
\]
In particular, after the finite initial range has been fixed, the \(L_2\)-transport object is detected by finitely many scalar Casimir traces.
\end{corollary}

\begin{proof}
The theorem gives a fixed invertible linear transformation between the multiplicity vector and the moment vector at each degree \(n\).
Hence the \(k[\sigma]\)-submodules generated by these two sequences are canonically isomorphic.
\end{proof}

\section{The PBW rank-one discriminant comparison}

The algebra of \S\S2--3 becomes a recurrence module for a chiral algebra
only after Hypothesis~\ref{hyp:finite-character-recurrence} supplies a
minimal shift polynomial.
The rank-one affine calculation supplies a different exact object:
the two-sheeted discriminant cover of the scalar bar series.
Its reversed branch divisor is the polynomial from which the divisor-core module
\(M(p)\) is formed.

\subsection{Two sources of holonomic constraint}

Let \(\mathfrak g=\mathfrak{sl}_2\), with exponent \(m_1=1\) and invariant
degree \(d_1=2\), and let \(A=\widehat{\mathfrak{sl}}_{2,k}\) at
non-critical level \(k\ne -2\).
The bar chain groups
\[
\Bbar^n(A)=\mathfrak g^{\otimes n}\otimes
\Omega^{n-1}(\operatorname{Conf}_n)
\]
carry two independent structural inputs:

\begin{enumerate}[label=\textup{(\Roman*)}]
\item \textbf{Invariant-theory input.}
The quadratic Casimir \(C_2\in Z(U\mathfrak g)\) acts on
\(\mathfrak g^{\otimes n}\) via the diagonal adjoint action,
and the bar differential is \(\mathfrak g\)-equivariant.
Therefore bar cohomology decomposes into isotypical components
\[
H^n(\Bbar(A))=\bigoplus_\lambda h_\lambda^{(n)}V(\lambda),
\]
and the dimension is
\(a_n=\sum_\lambda h_\lambda^{(n)}\dim V(\lambda)\).
The generating function for the invariant ring
\((\operatorname{Sym}^*\mathfrak{sl}_2^*)^{\mathfrak{sl}_2}\cong
k[C_2]\) is the Molien--Weyl series
\begin{equation}\label{eq:molien-weyl-g}
\sum_{p\ge 0}\dim(\operatorname{Sym}^p\mathfrak g)^{\mathfrak g}
\,t^p
=\frac{1}{1-t^2},
\end{equation}
a rational function with one invariant-theory pole.
This pole supplies one holonomic constraint on the bar cohomology
generating function through the PBW spectral sequence
(Theorem~\ref{thm:pbw-koszulness-criterion}):
the \(E_1\) page
\(E_1^{p,q}=H^q(\mathfrak g;\operatorname{Sym}^p
(\mathfrak g[-1]))\)
is controlled by the Molien--Weyl data, and the
Chevalley--Eilenberg vanishing \(H^1=H^2=0\) (Whitehead)
forces the spectral sequence to collapse early.

\item \textbf{Arnold input.}
The Orlik--Solomon dimension
\(\dim\Omega^{n-1}(\operatorname{Conf}_n)=(n-1)!\)
contributes a factorial-growth mode absent from the invariant
theory.
In the generating function, this manifests as the Euler-type
differential operator \(x\partial_x\); in the Riordan normalization below
it is the source of the polynomial coefficient in
\[
(n{+}1)a_n=(n{-}1)(2a_{n-1}+3a_{n-2}).
\]
\end{enumerate}

\begin{theorem}[PBW discriminant theorem for
 \texorpdfstring{$\widehat{\mathfrak{sl}}_2$}{sl2-hat}]
\label{thm:pbw-recurrence}
\index{PBW recurrence theorem}
Let \(A=\widehat{\mathfrak{sl}}_{2,k}\) at non-critical level.
Use the Riordan normalization
\[
P_{\mathfrak{sl}_2}(x)=\sum_{n\ge 0}a_nx^n,\qquad a_0=1.
\]
Then:
\begin{enumerate}[label=\textup{(\alph*)}]
\item
The scalar bar series satisfies
\[
x(1{+}x)P_{\mathfrak{sl}_2}(x)^2-(1{+}x)P_{\mathfrak{sl}_2}(x)+1=0.
\]
\item
The two-sheeted algebraic cover defined by this equation has branch
divisor
\[
\Delta_{\widehat{\mathfrak{sl}}_2}(x)=(1{+}x)(1{-}3x).
\]
\item
The reversed discriminant is
\[
p_{\widehat{\mathfrak{sl}}_2}(t)
=t^2\Delta_{\widehat{\mathfrak{sl}}_2}(t^{-1})
=(t-3)(t+1).
\]
\item
The divisor-core transport object attached to the scalar branch divisor is
\[
M\bigl(p_{\widehat{\mathfrak{sl}}_2}\bigr)
=k[t]/\bigl((t-3)(t+1)\bigr).
\]
It agrees with the intrinsic Casimir recurrence module \(\mathcal R_A\)
exactly when Hypothesis~\ref{hyp:finite-character-recurrence} holds for
this \(A\) with minimal polynomial \(p_{\widehat{\mathfrak{sl}}_2}\).
\end{enumerate}
\end{theorem}

\begin{remark}[Scope of the PBW comparison]
For every simple \(\mathfrak g\), once a discriminant polynomial or a
minimal recurrence polynomial has been supplied, the divisor-core
calculus of \S2 applies verbatim.
The rank-one theorem supplies the polynomial from the scalar branch cover;
it does not convert a P-recursive coefficient sequence into a constant-coefficient
shift recurrence.
\end{remark}

\subsection{Proof for \texorpdfstring{$\mathfrak{sl}_2$}{sl(2)}}
\label{subsec:pbw-recurrence-sl2}

The discriminant calculation for \(\mathfrak g=\mathfrak{sl}_2\)
(rank \(r=1\), exponent \(m_1=1\), degree \(d_1=2\)).

\medskip\noindent
\textbf{The Molien series.}
\(\operatorname{Sym}^p(\mathfrak{sl}_2)\cong
V(2p)\oplus V(2p-4)\oplus\cdots\),
so
\[
\dim(\operatorname{Sym}^p(\mathfrak{sl}_2))^{\mathfrak{sl}_2}
=\begin{cases}1&p\text{ even},\\0&p\text{ odd},\end{cases}
\]
and the Molien series is \(1/(1-t^2)\), with a single pole at
\(t^2=1\).

\medskip\noindent
\textbf{The PBW spectral sequence.}
The CE cohomology \(H^q(\mathfrak{sl}_2,M)=0\) for
\(q=1,2\) by Whitehead, and is nonzero only at \(q=0\) and
\(q=3\) (since \(\dim\mathfrak{sl}_2=3\)).
On the \(E_1\)~page,
\[
E_1^{p,0}=(\operatorname{Sym}^p\mathfrak{sl}_2)^{\mathfrak{sl}_2},
\qquad
E_1^{p,3}=(\operatorname{Sym}^p\mathfrak{sl}_2)_{\mathfrak{sl}_2},
\]
both nonzero only when \(p\) is even (one-dimensional).
Thus the \(E_1\)~page has exactly two nonzero rows,
each with one-dimensional entries at even columns.
The bar cohomology \(H^n\) is assembled from these two rows.

\medskip\noindent
\textbf{The algebraic branch equation.}
The Riordan generating function
\(P(x)=\sum_{n\ge 0}a_n x^n\) (where \(a_n\) are the Riordan
numbers) satisfies the algebraic equation
\[
x(1{+}x)P^2-(1{+}x)P+1=0
\]
of degree \(2\) over \(\mathbb Q(x)\),
with discriminant
\[
(1{+}x)^2-4x(1{+}x)=(1{+}x)(1{-}3x)
=\Delta_{\widehat{\mathfrak{sl}}_2}(x).
\]
The branch points are \(x=1/3\) and \(x=-1\); their reciprocals are
the roots \(3\) and \(-1\) of \(p_{\widehat{\mathfrak{sl}}_2}(t)\).

\medskip\noindent
\textbf{Intrinsic origin.}
This two-sheeted branch structure is compatible with two sources:
\begin{itemize}
\item
The Molien pole at \(t^2=1\) provides the invariant-theory
constraint: the even/odd periodicity of
\(E_1^{p,0}\) couples consecutive bar degrees through a
single Casimir conservation law.
\item
The Arnold factorial \((n{-}1)!\) provides the Euler constraint:
the P-recursive recurrence
\((n{+}1)a_n=(n{-}1)(2a_{n{-}1}+3a_{n{-}2})\)
has polynomial coefficients of degree 1 in \(n\),
matching the two-dimensional branch cover.
\end{itemize}
The discriminant module records the branch roots \(3\) and \(-1\).
It is not the same assertion as a constant-coefficient recurrence for the
coefficient sequence.

\subsection{Growth-mode factorization}

In rank~\(1\), the eigenvalue structure of
\(\Delta_{\widehat{\mathfrak{sl}}_2}\) admits a canonical
factorization.

\begin{theorem}[Growth-mode factorization for
 \texorpdfstring{$\widehat{\mathfrak{sl}}_2$}{sl2-hat}]
\label{thm:growth-mode-factorization}
\index{growth-mode factorization}
The spectral discriminant of \(\widehat{\mathfrak{sl}}_2\) factors as
\begin{equation}\label{eq:growth-mode-factor}
\Delta_{\widehat{\mathfrak{sl}}_2}(x)
=(1-3x)(1+x)
=(1-(\dim\mathfrak{sl}_2)\,x)\cdot
\Delta^{\mathrm{DS}}_{\mathfrak{sl}_2}(x),
\end{equation}
where:
\begin{enumerate}[label=\textup{(\roman*)}]
\item
\(\Delta^{\mathrm{DS}}_{\mathfrak{sl}_2}(x)=1+x\) has degree~\(1\).
\item
\(\Delta^{\mathrm{DS}}_{\mathfrak{sl}_2}\) divides
\(\Delta_{\mathrm{Vir}}(x)\).
\item
The dominant eigenvalue \(\dim\mathfrak{sl}_2=3\) is the
chain-group growth rate
\(\dim\mathfrak{sl}_2^{\otimes n}=3^n\).
\item
The remaining eigenvalue \(-1\)
\textup{(}the root of
\((\Delta^{\mathrm{DS}}_{\mathfrak{sl}_2})^{\rev}\)\textup{)}
encodes the
representation-theoretic cancellations in bar cohomology.
\end{enumerate}
\end{theorem}

\begin{proof}[Verification for \(\mathfrak{sl}_2\)]
\(\Delta_{\widehat{\mathfrak{sl}}_2}=(1{-}3x)(1{+}x)\).
Factor: \((1{-}3x)\) with \(3=\dim\mathfrak{sl}_2\) and
\(\Delta^{\mathrm{DS}}=(1{+}x)\), degree \(1=r\).
The Virasoro divisor is \(\Delta_{\mathrm{Vir}}=(1{-}3x)(1{+}x)\),
so \((1{+}x)\mid\Delta_{\mathrm{Vir}}\).
\end{proof}

\begin{remark}[Higher-rank input]
For a higher-rank simple \(\mathfrak g\), the statements below use only
the displayed recurrence polynomial.
If a separate computation supplies a factorization of
\(\Delta_{\widehat{\mathfrak g}}\), the divisor-core calculus compares
the factors without asserting full operator similarity.
\end{remark}

\begin{remark}[What the factorization means]
The dominant factor \((1-(\dim\mathfrak{sl}_2)\,x)\) is
\emph{visible on chain groups}: it is the radius of
convergence of
\(\sum\dim\mathfrak{sl}_2^{\otimes n}x^n
=3x/(1{-}3x)\),
requiring no bar differential.
The factor \(\Delta^{\mathrm{DS}}_{\mathfrak{sl}_2}\) records the
non-dominant branch root of the scalar series and is the common factor
with the Virasoro discriminant.
\end{remark}

\begin{corollary}[Divisor-core interpretation of the
 factorization]
\label{cor:growth-ds-divisor-cores}
Under the \(\widehat{\mathfrak{sl}}_2\) growth-mode factorization,
the discriminant-core module has a canonical exact sequence
\[
0\longrightarrow
\underbrace{\mathcal C_{t+1}
(p_{\widehat{\mathfrak{sl}}_2})}_{\text{DS-invariant core, dim }1}
\longrightarrow
M(p_{\widehat{\mathfrak{sl}}_2})
\longrightarrow
\underbrace{k[t]/(t{-}3)}_{\text{growth mode, dim }1}
\longrightarrow 0.
\]
Because
\(\gcd(t{-}3,\,
t{+}1)=1\), this sequence splits, giving an idempotent
projector onto each sector
(Proposition~\ref{prop:coprime-functional-calculus}).
Under Hypothesis~\ref{hyp:finite-character-recurrence} with
\(\mathcal R_{\widehat{\mathfrak{sl}}_2}\cong M(p_{\widehat{\mathfrak{sl}}_2})\),
the same sequence is the divisor-core decomposition of the intrinsic
Casimir recurrence module.
\end{corollary}

\begin{proof}
Apply Theorem~\ref{thm:divisor-core-calculus}(ii) and~(iv)
to the factorization
\(p_{\widehat{\mathfrak{sl}}_2}=(t{-}3)(t{+}1)\).
\end{proof}

\begin{remark}[Rank-one intrinsic comparison]
\label{rem:intrinsic-resolution}
Theorem~\ref{thm:pbw-recurrence} constructs the rank-one branch
polynomial
\[
\Delta=(1{-}3x)(1{+}x)
\]
from the PBW spectral sequence, the \(\mathfrak{sl}_2\) Casimir, and
the Arnold factor.
The intrinsic Kodaira--Spencer operator on a reduced cohomology space
requires the additional lift datum isolated in
Definition~\ref{def:geometric-divisor-core-lift-datum}.
\end{remark}

\section{Reduced spectral sectors and common transport}

\subsection{Divisors of the discriminant become subquotients}

Let \(A\) satisfy Hypothesis~\ref{hyp:finite-character-recurrence}, and write
\[
p_A(t)=t^{d_A}\Delta_A(t^{-1}).
\]

\begin{definition}[Reduced spectral sector]
Let \(S(x)\mid \Delta_A(x)\) be a divisor normalized by \(S(0)=1\), and define
\[
s(t):=S^{\rev}(t).
\]
The \emph{algebraic reduced spectral sector} of type \(S\) is the canonical subquotient
\[
H^{1,\alg}_S(A):=\bigl(\mathcal C_s(p_A)\bigr)^\vee.
\]
Equivalently,
\[
H^{1,\alg}_S(A)\cong H^{1,\alg}_{\red}(A)\Big/\bigl(\mathcal C_s(p_A)\bigr)^\perp.
\]
\end{definition}

\begin{theorem}[Characteristic polynomial of a sector]
\label{thm:sector-determinant}
For every divisor \(S\mid \Delta_A\) normalized by \(S(0)=1\),
\[
\bigl(H^{1,\alg}_S(A)\bigr)^\vee\cong k[t]/(s),
\]
and the induced endomorphism \(KS_{A,S}\) on \(H^{1,\alg}_S(A)\) satisfies
\[
\det(1-x\,KS_{A,S})=S(x).
\]
\end{theorem}

\begin{proof}
By Theorem~\ref{thm:divisor-core-calculus}\textup{(i)},
\[
\mathcal C_s(p_A)\cong k[t]/(s).
\]
The dual space therefore has dimension \(\deg s=\deg S\), and the endomorphism induced by \(\sigma^\vee\) has characteristic polynomial \(S\), exactly as in Theorem~\ref{thm:minimal-intrinsic-realization}.
\end{proof}

\begin{remark}
This is the fundamental passage from polynomial factorization to transport geometry:
every divisor of the discriminant defines a canonical finite-dimensional transport sector.
\end{remark}

\subsection{Common cores between two algebras}

Let \(A,B\) satisfy Hypothesis~\ref{hyp:finite-character-recurrence}.
Write
\[
p_A(t)=t^{d_A}\Delta_A(t^{-1}),
\qquad
p_B(t)=t^{d_B}\Delta_B(t^{-1}),
\qquad
g_{A,B}(t):=\gcd(p_A(t),p_B(t)).
\]

\begin{definition}[Maximal common transport core]
The \emph{maximal common transport core} of \(A\) and \(B\) is
\[
\mathcal C^{\max}_{A,B}:=k[t]/(g_{A,B}(t)).
\]
The corresponding reduced cohomological sector is the canonical subquotient
\[
H^{1,\alg}_{A\wedge B}:=\bigl(\mathcal C^{\max}_{A,B}\bigr)^\vee.
\]
\end{definition}

\begin{theorem}[Common-core exact sequences]
\label{thm:common-core-exact-sequences}
There are canonical short exact sequences
\[
0\longrightarrow \mathcal C^{\max}_{A,B}\longrightarrow \mathcal R_A\longrightarrow k[t]/(p_A/g_{A,B})\longrightarrow 0,
\]
\[
0\longrightarrow \mathcal C^{\max}_{A,B}\longrightarrow \mathcal R_B\longrightarrow k[t]/(p_B/g_{A,B})\longrightarrow 0.
\]
Dually, there are canonical quotient maps
\[
H^{1,\alg}_{\red}(A)\twoheadrightarrow H^{1,\alg}_{A\wedge B},
\qquad
H^{1,\alg}_{\red}(B)\twoheadrightarrow H^{1,\alg}_{A\wedge B}.
\]
\end{theorem}

\begin{proof}
The exact sequences on \(\mathcal R_A\) and \(\mathcal R_B\) are just Corollary~\ref{cor:spectral-defects}
applied to \(p_A\) and \(p_B\).
Dualizing yields the quotient maps.
\end{proof}

\begin{proposition}[Shift-compatible transport factors through the common core]
\label{prop:transport-factors-through-common-core}
Let
\[
F\colon \mathcal R_A\to \mathcal R_B
\]
be any \(k[t]\)-linear map.
Then \(F\) factors uniquely through the maximal common core:
\[
\mathcal R_A\twoheadrightarrow \mathcal C^{\max}_{A,B}\xrightarrow{\,m_u\,}\mathcal C^{\max}_{A,B}\hookrightarrow \mathcal R_B
\]
for a unique \(u\in \mathcal C^{\max}_{A,B}\).
Equivalently, every shift-compatible transport between the algebraic reduced sectors
\[
H^{1,\alg}_{\red}(B)\to H^{1,\alg}_{\red}(A)
\]
is governed by multiplication on the maximal common core.
\end{proposition}

\begin{proof}
Combine Theorem~\ref{thm:minimal-intrinsic-realization} with
Theorem~\ref{thm:factorization-through-common-core}.
\end{proof}

\begin{corollary}[Gcd bound for shift-compatible transport]
\label{cor:gcd-bound-transport}
Let \(A,B\) satisfy Hypothesis~\ref{hyp:finite-character-recurrence}.
Every \(k[t]\)-linear transport map \(\mathcal R_A\to \mathcal R_B\) is
determined by a unique element of \(k[t]/(g_{A,B})\).
In particular,
\[
\dim_k\Hom_{k[t]}(\mathcal R_A,\mathcal R_B)=\deg g_{A,B},
\]
and no nonzero shift-compatible transport exists when \(g_{A,B}=1\).
\end{corollary}

\begin{proof}
This is Theorem~\ref{thm:hom-equals-gcd} applied to
\(\mathcal R_A\cong M(p_A)\) and \(\mathcal R_B\cong M(p_B)\).
\end{proof}

\subsection{Squarefree common sectors}

\begin{proposition}[Squarefree common-sector rigidity]
\label{prop:squarefree-common-sector}
Suppose \(S(x)\) is a common divisor of \(\Delta_A(x)\) and \(\Delta_B(x)\), normalized by \(S(0)=1\), and assume \(S\) is squarefree.
Then after extension to a splitting field \(K\), the \(S\)-sectors of \(A\) and \(B\) are canonically isomorphic to the same squarefree spectral sheet:
\[
K\otimes_k \bigl(H^{1,\alg}_S(A)\bigr)^\vee
\cong \operatorname{Sh}(S)
\cong K\otimes_k \bigl(H^{1,\alg}_S(B)\bigr)^\vee.
\]
In particular, the two induced operators are conjugate as \(K[t]\)-modules.
\end{proposition}

\begin{proof}
By Theorem~\ref{thm:sector-determinant}, both dual sectors are isomorphic to \(k[t]/(s)\), where \(s=S^{\rev}\).
Apply Theorem~\ref{thm:squarefree-sheet}.
\end{proof}

\begin{remark}
On the squarefree locus, the discriminant already determines the full semisimple transport sector.
All genuinely new spectral geometry is therefore forced onto the repeated-root locus, or onto the higher layers of the characteristic hierarchy.
\end{remark}

\section{Examples}

\subsection{\(\mathfrak{sl}_2\), Virasoro, and the common two-sheet}

For the \(\mathfrak{sl}_2\) Drinfeld--Sokolov orbit one has
\[
\Delta_{\widehat{\mathfrak{sl}}_2}(x)
=
\Delta_{\mathrm{Vir}}(x)
=
(1-3x)(1+x),
\]
hence
\[
p_{\widehat{\mathfrak{sl}}_2}(t)=p_{\mathrm{Vir}}(t)=(t-3)(t+1)=t^2-2t-3.
\]
Thus the corresponding discriminant-core modules satisfy
\[
M(p_{\widehat{\mathfrak{sl}}_2})\cong M(p_{\mathrm{Vir}})
\cong k[t]/(t^2-2t-3),
\]
and the maximal common core is the entire transport module.
Because the polynomial is squarefree, the corresponding transport is already semisimple:
\[
\operatorname{Spec}(T)=\{3,-1\}.
\]

\begin{remark}
This is why the rank-one case is deceptive.
Every divisor is automatically coprime to its complement, so the common spectral sector splits completely.
Nothing in rank one foreshadows the Jordan and defect phenomena that appear as soon as a proper common factor survives but the whole operator does not.
\end{remark}

\subsection{\(\widehat{\mathfrak{sl}}_3\) versus \(W_3\): quadratic common core}

Assume the denominator chart
\[
\Delta_{\widehat{\mathfrak{sl}}_3}(x)=(1-8x)(1-3x-x^2),
\qquad
\Delta_{W_3}(x)=(1-x)(1-3x-x^2).
\]
Then
\[
p_{\widehat{\mathfrak{sl}}_3}(t)=(t-8)(t^2-3t-1),
\qquad
p_{W_3}(t)=(t-1)(t^2-3t-1).
\]

\begin{definition}[The \(\mathfrak{sl}_3/W_3\) common core]
Set
\[
g_{\DS}(t):=t^2-3t-1.
\]
The corresponding common sector is
\[
\mathcal C^{\max}_{\widehat{\mathfrak{sl}}_3,W_3}
\cong k[t]/(t^2-3t-1).
\]
\end{definition}

\begin{proposition}[Exact defect decomposition for \(\widehat{\mathfrak{sl}}_3/W_3\)]
\label{prop:sl3-w3-defect}
Under the displayed denominator chart, one has canonical exact sequences
\[
0\longrightarrow k[t]/(t^2-3t-1)\longrightarrow
k[t]/\bigl((t-8)(t^2-3t-1)\bigr)\longrightarrow k[t]/(t-8)\longrightarrow 0,
\]
\[
0\longrightarrow k[t]/(t^2-3t-1)\longrightarrow
k[t]/\bigl((t-1)(t^2-3t-1)\bigr)\longrightarrow k[t]/(t-1)\longrightarrow 0.
\]
Thus the non-common factors are the one-dimensional defect modules
\[
k[t]/(t-8)
\qquad\text{and}\qquad
k[t]/(t-1).
\]
\end{proposition}

\begin{proof}
This is Theorem~\ref{thm:common-core-exact-sequences} applied to the two recurrence polynomials.
\end{proof}

\begin{corollary}[All transport maps factor through the quadratic core]
\[
\Hom_{k[t]}\!\left(
k[t]/\bigl((t-8)(t^2-3t-1)\bigr),
k[t]/\bigl((t-1)(t^2-3t-1)\bigr)
\right)
\cong
k[t]/(t^2-3t-1).
\]
In particular, every shift-compatible transport between the two cyclic modules is controlled by the two-dimensional common core, and not by the one-dimensional defects.
\end{corollary}

\begin{proof}
Apply Theorem~\ref{thm:hom-equals-gcd}.
\end{proof}

\begin{proposition}[Explicit coprime-locus projectors]
\label{prop:explicit-sl3-projectors}
Because
\[
\gcd(t^2-3t-1,t-8)=\gcd(t^2-3t-1,t-1)=1,
\]
the common quadratic sector splits idempotently on both sides.

\begin{enumerate}[label=\textup{(\roman*)}]
\item For \(\widehat{\mathfrak{sl}}_3\), a Bezout identity
\[
\frac1{39}(t^2-3t-1)-\frac{t+5}{39}(t-8)=1
\]
gives the projector
\[
P^{\mathrm{core}}_{\widehat{\mathfrak{sl}}_3}
=
-\frac{(T+5I)(T-8I)}{39}.
\]

\item For \(W_3\), a Bezout identity
\[
-\frac13(t^2-3t-1)+\frac{t-2}{3}(t-1)=1
\]
gives the projector
\[
P^{\mathrm{core}}_{W_3}
=
\frac{(T-2I)(T-I)}{3}.
\]
\end{enumerate}

Each image has characteristic polynomial \(t^2-3t-1\), hence eigenvalues
\[
\frac{3+\sqrt{13}}2,\qquad \frac{3-\sqrt{13}}2.
\]
\end{proposition}

\begin{proof}
Apply Proposition~\ref{prop:coprime-functional-calculus} with \(g=t^2-3t-1\) and \(p/g=t-8\) or \(t-1\), respectively.
\end{proof}

\begin{remark}
On this denominator chart the common quadratic factor is rigid, while the
dominant eigenvalues \(8\) and \(1\) lie in different one-dimensional
defect modules.
The preserved shift-compatible object is the quadratic common core.
\end{remark}

\section{Geometric lift data}

\subsection{The four exact distinctions}

The preceding algebra enforces four exact distinctions.

\begin{enumerate}[label=\textup{(\arabic*)}]
\item The discriminant is \emph{not} the final object; it is the characteristic polynomial of the final object.

\item A divisor of the discriminant does \emph{not} in general define a projector; it defines a canonical subquotient.

\item Repeated roots are \emph{not} a pathology to be ignored; they are the source of genuine indecomposable transport sectors.

\item Drinfeld--Sokolov transport is \emph{not} full operator similarity in rank \(\ge 2\); the corrected invariant is the common divisor-core.
\end{enumerate}

\subsection{Characteristic hierarchy comparison}

Work on a Lie-theoretic locus where the independent characteristic data reduce to the effective quadruple
\[
(\Gamma_A,\Delta_A,\Pi_A,H_A),
\]
with \(\Theta_A\) recoverable from \(\Gamma_A\). On the
uniform-weight lane (Definition~\ref{def:scalar-lane}),
\(\Gamma_A = \kappaChHodge(A)\Lambda\). Without that scalar-lane hypothesis,
the equality \(\Gamma_A=\kappaChHodge(A)\Lambda\) is not part of the datum. The genuinely
non-scalar information begins at the spectral level \(\Delta_A\).
The construction above refines that spectral level from a polynomial to an actual module.

\begin{proposition}[The \(L_1\)--\(L_2\) package on the one-channel squarefree locus]
Assume:
\begin{enumerate}[label=\textup{(\roman*)}]
\item \(A\) and \(B\) lie on a scalar-saturated one-channel locus, so their universal Maurer--Cartan classes are determined by \(\kappaChHodge(A)\) and \(\kappaChHodge(B)\);
\item \(\Delta_A=\Delta_B\);
\item \(\Delta_A\) is squarefree.
\end{enumerate}
Then the \(L_2\)-transport sectors of \(A\) and \(B\) are conjugate as
\(k[t]\)-modules after extension to a splitting field, and the \(L_1\)-data
coincide if and only if \(\kappaChHodge(A)=\kappaChHodge(B)\).
Thus, on this locus, the realized \(L_1\)--\(L_2\) package is determined by the pair
\[
(\kappa,\Delta).
\]
\end{proposition}

\begin{proof}
The squarefree part follows from Proposition~\ref{prop:squarefree-common-sector}.
The \(\kappa\)-statement is exactly the scalar-saturated one-channel hypothesis.
\end{proof}

\begin{remark}
If \(\Theta\) is rigidly determined by \(\kappa\), then beyond the scalar line there are two remaining directions:
the Jordan geometry of repeated roots at \(L_2\), and the periodic or higher-channel data above it.
\end{remark}

\subsection{Geometric divisor-core lift data}

\begin{definition}[Geometric divisor-core lift datum]
\label{def:geometric-divisor-core-lift-datum}
Let \(A\) satisfy Hypothesis~\ref{hyp:finite-character-recurrence}.
A \emph{geometric divisor-core lift datum} for \(A\) consists of a finite-dimensional
subquotient
\[
H^1_{\red}(A)
\]
of \(H^1(\Bbar(A))\), together with an endomorphism
\[
KS_A\in \End\bigl(H^1_{\red}(A)\bigr),
\]
and a \(k[t]\)-module isomorphism
\[
\eta_A\colon \bigl(H^1_{\red}(A)\bigr)^\vee\xrightarrow{\sim}\mathcal R_A
\]
which intertwines \(KS_A^\vee\) with the shift operator \(\sigma\) on \(\mathcal R_A\).
\end{definition}

\begin{theorem}[Divisor-core consequences of a geometric lift datum]
\label{thm:geometric-lift-datum-consequences}
Let \(A\) satisfy Hypothesis~\ref{hyp:finite-character-recurrence}, and fix a
geometric divisor-core lift datum in the sense of
Definition~\ref{def:geometric-divisor-core-lift-datum}.
Then:
\begin{enumerate}[label=\textup{(\roman*)}]
\item for every divisor \(S(x)\mid \Delta_A(x)\) normalized by \(S(0)=1\), writing \(s(t)=S^{\rev}(t)\), there is a canonical subquotient
\[
H^1_S(A)
\]
of \(H^1_{\red}(A)\) whose dual is the divisor core
\[
\bigl(H^1_S(A)\bigr)^\vee \cong \mathcal C_s(p_A);
\]

\item these sectors satisfy the divisor-lattice identities of Theorem~\ref{thm:divisor-core-calculus}, up to canonical isomorphism of subquotients;

\item the obstruction to an idempotent splitting of \(H^1_S(A)\) is the class
\[
\varepsilon_S(A):=\varepsilon_s(p_A)\in
\Ext^1_{k[t]}\!\bigl(k[t]/(p_A/s),k[t]/(s)\bigr)
\cong k[t]/\gcd(s,p_A/s),
\]
and this class vanishes if and only if
\[
\gcd\!\left(S,\frac{\Delta_A}{S}\right)=1.
\]
\end{enumerate}
\end{theorem}

\begin{proof}
The isomorphism \(\eta_A\) identifies \(\bigl(H^1_{\red}(A)\bigr)^\vee\)
with \(M(p_A)\).
For \(S\mid \Delta_A\), the submodule \(\mathcal C_s(p_A)\subset M(p_A)\)
therefore defines the quotient
\[
H^1_{\red}(A)\twoheadrightarrow H^1_S(A):=\bigl(\mathcal C_s(p_A)\bigr)^\vee.
\]
The divisor-lattice identities are dual to Theorem~\ref{thm:divisor-core-calculus}.
The obstruction class and its vanishing criterion are
Corollary~\ref{cor:repeated-roots-extension}, applied to \(g=s\) and
\(p=p_A\); reversal identifies \(\gcd(s,p_A/s)=1\) with
\(\gcd(S,\Delta_A/S)=1\).
\end{proof}

\begin{proposition}[Primary quotient filtration from a lift datum]
\label{prop:primary-quotient-filtration-lift}
Let \(A\) carry a geometric divisor-core lift datum, let \(q\) be an
irreducible factor of \(p_A\), and let \(1\le m\le \nu_q(p_A)\).
The primary divisor core
\[
\mathcal J_q^m(p_A)\subset \mathcal R_A
\]
defines a canonical quotient
\[
H^1_{\red}(A)\twoheadrightarrow
H^1_{q,m}(A):=\bigl(\mathcal J_q^m(p_A)\bigr)^\vee.
\]
The quotients form an inverse filtration dual to
\[
0\subset \mathcal J_q^1(p_A)\subset\cdots\subset \mathcal J_q^m(p_A).
\]
If \(q=t-\lambda\) splits, the induced operator \(KS_A-\lambda\) on
\(H^1_{q,m}(A)\) is nilpotent of order \(m\) and has one Jordan block of
size \(m\).
\end{proposition}

\begin{proof}
This is Theorem~\ref{thm:primary-jordan-filtration} transported through
the isomorphism \(\eta_A\) and then dualized.
\end{proof}

\begin{theorem}[Geometric common-core factorization]
\label{thm:geometric-common-core-factorization}
Let \(A\) and \(B\) carry geometric divisor-core lift data.
Let \(F\) be an exact factorization functor or reduction procedure carrying
\(A\) to \(B\), and assume that the induced map on the lift data is
shift-compatible, equivalently that its dual is a \(k[t]\)-linear map
\[
F^\vee_{\red}\colon \mathcal R_A\to \mathcal R_B.
\]
Then \(F^\vee_{\red}\) factors uniquely through the maximal common core
\[
\mathcal R_A\twoheadrightarrow \mathcal C^{\max}_{A,B}
\xrightarrow{\,m_u\,}\mathcal C^{\max}_{A,B}\hookrightarrow \mathcal R_B
\]
for a unique \(u\in \mathcal C^{\max}_{A,B}\).
Dually, the reduced geometric transport factors through
\[
H^{1}_{A\wedge B}.
\]
\end{theorem}

\begin{proof}
Apply Proposition~\ref{prop:transport-factors-through-common-core} to
the \(k[t]\)-linear map \(F^\vee_{\red}\), then dualize the resulting
factorization.
\end{proof}

\begin{theorem}[Drinfeld--Sokolov common-core transport under lift data]
\label{thm:geometric-ds-common-core}
Let
\[
A=\widehat{\mathfrak g}_k,
\qquad
W=W_k(\mathfrak g),
\qquad
S_{\mathfrak g,k}(x):=\gcd(\Delta_A(x),\Delta_W(x)),
\qquad
s_{\mathfrak g,k}(t):=S_{\mathfrak g,k}^{\rev}(t).
\]
The gcd in \(k[x]\) is normalized by \(S_{\mathfrak g,k}(0)=1\).
Assume \(A\) and \(W\) carry geometric divisor-core lift data and that
Drinfeld--Sokolov reduction induces a shift-compatible map on the corresponding
Casimir recurrence modules.
Then the induced transport factors through the reduced common-core sector
\[
H^1_{S_{\mathfrak g,k}}(-).
\]
For \(\mathfrak{sl}_2\), \(S_{\mathfrak g,k}=(1-3x)(1+x)\), so the common
core is the whole discriminant-core object.
On the \(\mathfrak{sl}_3/W_3\) denominator chart above,
\[
S_{\mathfrak g,k}(x)=1-3x-x^2,\qquad
s_{\mathfrak g,k}(t)=t^2-3t-1,
\]
and the two linear factors \(t-8\) and \(t-1\) are the defect modules.
\end{theorem}

\begin{proof}
The reversal of the common discriminant factor is
\(\gcd(p_A,p_W)=s_{\mathfrak g,k}\).
The factorization is Theorem~\ref{thm:geometric-common-core-factorization}
with this gcd.
The \(\mathfrak{sl}_2\) and \(\mathfrak{sl}_3/W_3\) statements are the
displayed factorizations in the two examples above.
\end{proof}

\section{Polynomial Shadows and Transport Geometry}

The spectral discriminant \(\Delta_A\) determines the cyclic module
\[
M(p_A),\qquad p_A(t)=t^{d_A}\Delta_A(t^{-1}).
\]
Divisors of \(\Delta_A\) determine divisor cores in \(M(p_A)\); gcds
determine shift-compatible maps; repeated roots determine the extension
classes \(\varepsilon_g(p_A)\); coprime decompositions determine idempotent
projectors; squarefree factors determine semisimple spectral sheets.

\begin{construction}[Divisor-level transport datum for DS reductions]
\label{constr:casimir-divisor-shadow}
Let
\[
V=V_k(\mathfrak g),\qquad W=W_k(\mathfrak g,f)
\]
be a Drinfeld--Sokolov pair for which discriminant polynomials
\(\Delta_V,\Delta_W\) and reversed polynomials \(p_V,p_W\) have been fixed.
The divisor-level transport datum is the following purely algebraic package.
\begin{enumerate}[label=\textup{(\roman*)}]
\item The cyclic transport object is
\[
M(p_V)=k[t]/(p_V),\qquad p_V(t)=t^{d_V}\Delta_V(t^{-1}),
\]
and its subquotients are organized by Theorem~\ref{thm:divisor-core-calculus}.

\item If \(\Delta_W\mid\Delta_V\), write
\[
G_{V/W}(x):=\Delta_V(x)/\Delta_W(x),
\qquad
g_{V/W}(t):=G_{V/W}^{\rev}(t).
\]
Then
\[
p_V(t)=p_W(t)\,g_{V/W}(t),
\]
and the canonical quotient
\[
M(p_V)\twoheadrightarrow M(p_W)
\]
has kernel
\[
\mathcal C_{g_{V/W}}(p_V)\cong k[t]/(g_{V/W}).
\]

\item Without divisibility, the maximal functorial comparison is still defined:
\[
g_{V,W}(t):=\gcd(p_V(t),p_W(t)),\qquad
\mathcal C^{\max}_{V,W}=k[t]/(g_{V,W}).
\]
Every shift-compatible map \(M(p_V)\to M(p_W)\) factors through this common core.

\item For a finite reduction graph whose edge \(\lambda\to\mu\) carries
polynomials \(p_\lambda,p_\mu\), an edge with \(p_\mu\mid p_\lambda\)
acts by the quotient \(M(p_\lambda)\twoheadrightarrow M(p_\mu)\).
Its kernel is the divisor core \(\mathcal C_{p_\lambda/p_\mu}(p_\lambda)\);
if divisibility fails, the edge comparison is the common-core factorization
through \(k[t]/\gcd(p_\lambda,p_\mu)\).
\end{enumerate}
\end{construction}
